\chapter{Modular Hamiltonians and the Replica Formula}
\label{app:modular}

This appendix records the entanglement formulae used in
Chapters~6, 11, and~12. It introduces no new physical hypothesis.
The working identity of the main text remains
Proposition~\ref{prop:ward}. What follows is the language in which
that identity was written.

\section{States, reductions, and von Neumann entropy}
\label{sec:appB-vN}

Let \(\mathcal{H}\) be the Hilbert space of the microscopic theory
and \(\rho\) a density operator on \(\mathcal{H}\). For a von
Neumann algebra \(\mathcal{A}_D\) of observables associated with a
region \(D\), the reduced state \(\rho_D\) is the restriction of
\(\rho\) to \(\mathcal{A}_D\). In a tensor-product approximation
\(\mathcal{H}\simeq\mathcal{H}_D\otimes\mathcal{H}_{D^c}\),
\begin{equation}
  \label{eq:appB-rhoD}
  \rho_D=\mathrm{Tr}_{\mathcal{H}_{D^c}}\rho.
\end{equation}
Continuum quantum field theory does not have that tensor-product
factorization. The algebraic definition (restriction to
\(\mathcal{A}_D\)) is the honest one; equation~\eqref{eq:appB-rhoD}
is the formula used in every explicit replica computation in the
book.

The von Neumann entropy and modular Hamiltonian of \(\rho_D\) are
\begin{equation}
  \label{eq:appB-SH}
  S_D=-\mathrm{Tr}(\rho_D\ln\rho_D),
  \qquad
  H_D=-\ln\rho_D,
\end{equation}
so that
\begin{equation}
  \label{eq:appB-S-eq-H}
  S_D=\langle H_D\rangle_{\rho_D}.
\end{equation}
If \(\rho_D\) has kernel, \(\ln\rho_D\) is defined on the support
of \(\rho_D\) and \(S_D\) may be infinite. In quantum field theory
\(S_D\) is ultraviolet divergent. Differences \(\delta S_D\) at
fixed cutoff, and mutual informations, are the finite objects used
in the main text.

\section{Relative entropy and mutual information}
\label{sec:appB-rel}

For two states \(\rho\) and \(\sigma\) on the same algebra,
\begin{equation}
  \label{eq:appB-rel}
  S(\rho\|\sigma)
  =\mathrm{Tr}(\rho\ln\rho)-\mathrm{Tr}(\rho\ln\sigma)
\end{equation}
is non-negative and vanishes if and only if \(\rho=\sigma\)
(when the expression is finite). Mutual information of two
regions \(A\) and \(B\) is the relative entropy of \(\rho_{AB}\)
with respect to \(\rho_A\otimes\rho_B\),
\begin{equation}
  \label{eq:appB-MI}
  I(A:B)
  =S(\rho_{AB}\|\rho_A\otimes\rho_B)
  =S(A)+S(B)-S(A\cup B),
\end{equation}
the second equality holding when the three von Neumann entropies
are regulated in the same scheme. Positivity of \(I(A:B)\) is
therefore positivity of relative entropy. That is the source of
the sign in Definition~\ref{def:Uent} and of
\(\Gamma(d)\ge 0\) in Chapter~14.

Strong subadditivity,
\begin{equation}
  \label{eq:appB-SSA}
  S(A\cup B)+S(B\cup C)
  \ge S(B)+S(A\cup B\cup C),
\end{equation}
is recorded for completeness. The main text does not use it
beyond the positivity of \(I(A:B)\).

\section{The replica formula}
\label{sec:appB-replica}

For integer \(n\ge 1\),
\begin{equation}
  \label{eq:appB-replica}
  \mathrm{Tr}\,(\rho_D)^n
  =Z_n/Z_1^n,
\end{equation}
where \(Z_n\) is the partition function of the microscopic theory
on an \(n\)-sheeted cover branched over the edge
\(\mathcal{B}=\partial D\cap\Sigma\). Analytic continuation in
\(n\) and the identity
\begin{equation}
  \label{eq:appB-replica-S}
  S_D
  =\lim_{n\to 1}
   \bigl(1-n\partial_n\bigr)\ln\mathrm{Tr}\,(\rho_D)^n
\end{equation}
produce the von Neumann entropy. Differentiating at \(n=1\)
inserts a single modular Hamiltonian: the first variation
\(\delta S_D=\delta\langle H_D\rangle\) used in Chapter~6 is the
replica formula at first order, not a second definition.

In Presentation~B the cover is a branched cover of the auxiliary
manifold, later identified with the emergent spacetime by
Proposition~\ref{prop:scaffold}. In Presentation~A the same
construction requires Problem~A of Chapter~17.

\section{Bisognano--Wichmann and small diamonds}
\label{sec:appB-BW}

On Minkowski space, let \(W\) be a Rindler wedge. The
Bisognano--Wichmann theorem identifies the modular Hamiltonian of
the vacuum restricted to \(W\) with the boost generator that
fixes the edge,
\begin{equation}
  \label{eq:appB-BW}
  H_W
  =\frac{2\pi}{\kappa}
   \int_{\Sigma\cap W}
   \mathrm{d}\Sigma^\mu\,\xi^\nu T_{\mu\nu},
\end{equation}
where \(\xi^\mu\) is the boost Killing vector and
\(\kappa\) its surface gravity (equal to \(1\) in standard Rindler
coordinates, restored here for matching~\eqref{eq:HD-boost}).

A small causal diamond in a slowly curved emergent geometry is,
to leading order in Riemann normal coordinates~\eqref{eq:RNC}, a
compact analogue of two opposing Rindler wedges glued at a
2-sphere. The working formula~\eqref{eq:HD-boost} is that analogue:
the same boost-energy integral plus an area term
\(H_{\mathrm{area}}\) whose first variation is
\(\delta A/(4\Geff)\) after Chapter~8 has named \(\Geff\). The
error is \(O(\ell^3)\) in the diamond radius, which is the order
dropped in Proposition~\ref{prop:ward}.

The identification is exact for free fields in Rindler space. It
is an operator-product approximation for interacting condensate
fluctuations. The book does not claim a Bisognano--Wichmann
theorem for the full matrix potential~\eqref{eq:V}.

\section{Fixed volume}
\label{sec:appB-volume}

Jacobson's entanglement-equilibrium condition is imposed at fixed
spatial volume of the maximal slice. In Riemann normal coordinates
the volume of a small ball of radius \(\ell\) is~\eqref{eq:vol-exp}.
A pure conformal variation \(\delta g_{\mu\nu}=\chi\eta_{\mu\nu}\)
changes that volume at order \(\ell^3\). Holding \(V\) fixed
removes that breathing from the comparison and is why the conformal
block \(\psi\) and the tensor block \(T_{\mu\nu}\) do not enter
the first law on the same footing.

Fixed volume is a constraint on the variations admitted in
Assumption~\ref{ass:EE}. It is not a statement that \(\psi\) is
massive. That statement is Assumption~\ref{ass:gap-psi}.

\section{Boundary entropy versus diamond entropy}
\label{sec:appB-two}

Chapter~12 uses a different functional: a local density
\(S_{\mathrm{ent}}(h_{ij},\phi)\) on \(\partial M\). The replica
construction of that density is the same formula,
\eqref{eq:appB-replica-S}, applied to a thin tubular neighbourhood
of the boundary and then localized as in
Assumption~\ref{ass:Sent-density}. The two functionals must not be
substituted for each other. Section~\ref{sec:ch12-diamonds}
already assigned them different jobs.

\section{What this appendix will not do}
\label{sec:appB-not}

It will not develop Tomita--Takesaki theory beyond the definition
\(H_D=-\ln\rho_D\). It will not compute \(S_D\) for the matrix
potential~\eqref{eq:V}. It will not prove Assumption~\ref{ass:EE}.
Those are Problem~B of Chapter~17 and the physical input of
Chapter~6, respectively.

\section{Exercises}
\label{sec:appB-exercises}

\begin{enumerate}
  \item \textbf{Replica derivative.}
    From~\eqref{eq:appB-replica-S} and
    \(\mathrm{Tr}\,\rho^n=\exp(\ln\mathrm{Tr}\,\rho^n)\), recover
    \(S=\langle H\rangle\) at \(n=1\).
  \item \textbf{Relative entropy and Klein's inequality.}
    For finite-dimensional density matrices, prove
    \(S(\rho\|\sigma)\ge 0\). Deduce \(I(A:B)\ge 0\).
  \item \textbf{Rindler temperature.}
    In standard Rindler coordinates, show that the Unruh
    temperature seen by an accelerated observer is
    \(\kappa/2\pi\). Match the factor in~\eqref{eq:appB-BW}.
  \item \textbf{Area divergence.}
    For a free scalar in Minkowski space, the unregulated
    \(S_D\) diverges as \(\mathrm{Area}(\mathcal{B})/\epsilon^2\).
    Why do mutual information and first variations \(\delta S_D\)
    at fixed cutoff remain useful?
  \item \textbf{Two functionals.}
    Write four sentences explaining why a small-diamond \(S_D\)
    cannot be used as the boundary density of Definition~\ref{def:HBEV}.
\end{enumerate}